\documentclass[10pt]{beamer}
\usetheme{metropolis}
\usepackage{graphicx}
\usepackage{listings}
\usepackage{booktabs}
\usepackage{xcolor}
\usepackage{hyperref}

\definecolor{codebg}{HTML}{F5F5F5}
\definecolor{codegreen}{HTML}{2E7D32}
\definecolor{codegray}{HTML}{757575}
\definecolor{codeblue}{HTML}{1565C0}

\lstdefinestyle{rstyle}{
  language=R, backgroundcolor=\color{codebg},
  basicstyle=\ttfamily\scriptsize,
  keywordstyle=\color{codeblue}\bfseries,
  stringstyle=\color{codegreen},
  commentstyle=\color{codegray}\itshape,
  breaklines=true, frame=single,
  rulecolor=\color{codegray!30}, showstringspaces=false, tabsize=2,
  morekeywords={library,ggplot,aes,geom_col,geom_histogram,geom_point,
    geom_boxplot,geom_line,geom_bar,geom_text,coord_flip,scale_x_log10,
    theme_minimal,theme,element_text,element_blank,labs,
    fct_reorder,count,slice_max,mutate,filter,ggsave,position_stack}
}
\lstdefinestyle{pythonstyle}{
  language=Python, backgroundcolor=\color{codebg},
  basicstyle=\ttfamily\scriptsize,
  keywordstyle=\color{codeblue}\bfseries,
  stringstyle=\color{codegreen},
  commentstyle=\color{codegray}\itshape,
  breaklines=true, frame=single,
  rulecolor=\color{codegray!30}, showstringspaces=false, tabsize=4,
  morekeywords={import,from,as,pd,plt,sns,np,fig,ax,subplots,savefig}
}

\title{Comparative Lab: First Plots in R vs Python}
\subtitle{Workshop 1 --- Module 10}
\author{Prof.Asoc.\ Endri Raco}
\institute{Data Visualization for Data Scientists \\ UNYT Tirana}
\date{2025/26}

\begin{document}

\begin{frame}
  \titlepage
\end{frame}

% ================================================================
\begin{frame}{Module Roadmap}
  \begin{enumerate}
    \item Load the Google Play Store dataset in R and Python
    \item Produce six chart types side by side: bar, histogram, scatter,
          boxplot, line, stacked bar
    \item Apply Modules 1--9 principles to every chart
    \item Build a six-panel EDA dashboard in both languages
  \end{enumerate}
  \begin{exampleblock}{Learning Outcome}
    You will produce identical, publication-quality charts in both R and
    Python, establishing the comparative fluency that continues through
    all nine workshops.
  \end{exampleblock}
\end{frame}

% ================================================================
\section{Data Loading}
% ================================================================

\begin{frame}[fragile]{Loading the Dataset: R vs Python}
  \begin{columns}[T]
    \begin{column}{0.48\textwidth}
      \textbf{R (tidyverse)}
\begin{lstlisting}[style=rstyle]
library(tidyverse)

apps <- read_csv(
  "googleplaystore.csv")

# Quick inspection
dim(apps)       # 10841 x 13
glimpse(apps)
summary(apps$Rating)

# Clean
apps_clean <- apps |>
  filter(Type %in%
    c("Free", "Paid")) |>
  mutate(
    Reviews = as.numeric(Reviews))
\end{lstlisting}
    \end{column}
    \begin{column}{0.48\textwidth}
      \textbf{Python (pandas)}
\begin{lstlisting}[style=pythonstyle]
import pandas as pd
import numpy as np

apps = pd.read_csv(
    "googleplaystore.csv")

# Quick inspection
print(apps.shape) # (10841, 13)
print(apps.info())
print(apps["Rating"].describe())

# Clean
apps_clean = (apps
    .query("Type in ['Free','Paid']")
    .assign(Reviews=lambda d:
        pd.to_numeric(
            d["Reviews"],
            errors="coerce")))
\end{lstlisting}
    \end{column}
  \end{columns}
\end{frame}

% ================================================================
\section{Chart 1: Horizontal Bar}
% ================================================================

\begin{frame}[fragile]{Bar Chart --- R (ggplot2)}
  \begin{columns}[T]
    \begin{column}{0.52\textwidth}
\begin{lstlisting}[style=rstyle]
apps_clean |>
  count(Category, sort = TRUE) |>
  slice_max(n, n = 10) |>
  mutate(Category =
    fct_reorder(Category, n)) |>
  ggplot(aes(x = Category,
             y = n)) +
  geom_col(fill = "#1565C0",
           width = 0.6) +
  geom_text(aes(label =
    scales::comma(n)),
    hjust = -0.15, size = 3,
    fontface = "bold") +
  coord_flip() +
  theme_minimal() +
  theme(panel.grid.major.y =
    element_blank()) +
  labs(x = NULL, y = "Count",
    title = "Top 10 Categories")
\end{lstlisting}
    \end{column}
    \begin{column}{0.45\textwidth}
      \includegraphics[width=\textwidth]{figures_m10/bar_r}
    \end{column}
  \end{columns}
\end{frame}

% ================================================================
\begin{frame}[fragile]{Bar Chart --- Python (matplotlib)}
  \begin{columns}[T]
    \begin{column}{0.52\textwidth}
\begin{lstlisting}[style=pythonstyle]
top10 = (apps_clean
    .value_counts("Category")
    .head(10)
    .sort_values())

fig, ax = plt.subplots(
    figsize=(6, 5))
ax.barh(top10.index, top10.values,
    color="#2E7D32", height=0.6)

for i, v in enumerate(top10.values):
    ax.text(v + 15, i, f"{v:,}",
        va="center", fontsize=7,
        fontweight="bold")

for sp in ["top","right","left"]:
    ax.spines[sp].set_visible(False)
ax.set_xticks([])
ax.tick_params(left=False)
ax.set_title("Top 10 Categories",
    fontweight="bold")
plt.tight_layout()
fig.savefig("bar.pdf")
\end{lstlisting}
    \end{column}
    \begin{column}{0.45\textwidth}
      \includegraphics[width=\textwidth]{figures_m10/bar_py}
    \end{column}
  \end{columns}
\end{frame}

% ================================================================
\section{Chart 2: Histogram}
% ================================================================

\begin{frame}[fragile]{Histogram --- R (ggplot2)}
  \begin{columns}[T]
    \begin{column}{0.52\textwidth}
\begin{lstlisting}[style=rstyle]
mean_r <- mean(apps_clean$Rating,
               na.rm = TRUE)

ggplot(apps_clean,
       aes(x = Rating)) +
  geom_histogram(
    bins = 40,
    fill = "#1565C0",
    color = "white",
    linewidth = 0.3) +
  geom_vline(
    xintercept = mean_r,
    color = "#E53935",
    linewidth = 1,
    linetype = "dashed") +
  annotate("text",
    x = mean_r - 0.3,
    y = Inf, vjust = 1.5,
    label = paste0("Mean: ",
      round(mean_r, 2)),
    fontface = "bold",
    color = "#E53935", size = 3) +
  theme_minimal() +
  labs(title = "Rating Distribution",
       y = "Count")
\end{lstlisting}
    \end{column}
    \begin{column}{0.45\textwidth}
      \includegraphics[width=\textwidth]{figures_m10/hist_r}
    \end{column}
  \end{columns}
\end{frame}

% ================================================================
\begin{frame}[fragile]{Histogram --- Python (matplotlib)}
  \begin{columns}[T]
    \begin{column}{0.52\textwidth}
\begin{lstlisting}[style=pythonstyle]
ratings = apps_clean["Rating"].dropna()
mean_r = ratings.mean()

fig, ax = plt.subplots(figsize=(6,4))
ax.hist(ratings, bins=40,
    color="#2E7D32",
    edgecolor="white", lw=0.3)

ax.axvline(x=mean_r,
    color="#E53935", lw=2,
    linestyle="--")
ax.text(mean_r + 0.05,
    ax.get_ylim()[1] * 0.92,
    f"Mean: {mean_r:.2f}",
    fontsize=8, color="#E53935",
    fontweight="bold")

ax.set_xlabel("Rating")
ax.set_ylabel("Frequency")
ax.set_title("Rating Distribution",
    fontweight="bold")
ax.spines["top"].set_visible(False)
ax.spines["right"].set_visible(False)
\end{lstlisting}
    \end{column}
    \begin{column}{0.45\textwidth}
      \includegraphics[width=\textwidth]{figures_m10/hist_py}
    \end{column}
  \end{columns}
\end{frame}

% ================================================================
\section{Chart 3: Scatterplot}
% ================================================================

\begin{frame}[fragile]{Scatter --- R (ggplot2)}
  \begin{columns}[T]
    \begin{column}{0.52\textwidth}
\begin{lstlisting}[style=rstyle]
apps_clean |>
  filter(!is.na(Rating),
         Reviews > 0) |>
  ggplot(aes(x = Reviews,
             y = Rating,
             color = Type)) +
  geom_point(alpha = 0.3,
             size = 0.8) +
  scale_x_log10(
    labels = scales::comma) +
  scale_color_manual(
    values = c(Free = "#1565C0",
               Paid = "#E53935")) +
  theme_minimal() +
  labs(
    title = "Reviews vs Rating",
    x = "Reviews (log scale)",
    y = "Rating")
\end{lstlisting}
    \end{column}
    \begin{column}{0.45\textwidth}
      \includegraphics[width=\textwidth]{figures_m10/scatter_r}
    \end{column}
  \end{columns}
\end{frame}

% ================================================================
\begin{frame}[fragile]{Scatter --- Python (matplotlib)}
  \begin{columns}[T]
    \begin{column}{0.52\textwidth}
\begin{lstlisting}[style=pythonstyle]
df = apps_clean.dropna(
    subset=["Rating"])
df = df[df["Reviews"] > 0]

fig, ax = plt.subplots(figsize=(6,5))
palette = {"Free":"#2E7D32",
           "Paid":"#E53935"}

for t, c in palette.items():
    mask = df["Type"] == t
    ax.scatter(
        df.loc[mask, "Reviews"],
        df.loc[mask, "Rating"],
        s=8, c=c, alpha=0.3,
        edgecolors="none",
        label=t)

ax.set_xscale("log")
ax.set_xlabel("Reviews (log)")
ax.set_ylabel("Rating")
ax.set_title("Reviews vs Rating",
    fontweight="bold")
ax.legend(fontsize=7)
ax.spines["top"].set_visible(False)
ax.spines["right"].set_visible(False)
\end{lstlisting}
    \end{column}
    \begin{column}{0.45\textwidth}
      \includegraphics[width=\textwidth]{figures_m10/scatter_py}
    \end{column}
  \end{columns}
\end{frame}

% ================================================================
\section{Chart 4: Boxplot}
% ================================================================

\begin{frame}[fragile]{Boxplot --- R (ggplot2)}
  \begin{columns}[T]
    \begin{column}{0.52\textwidth}
\begin{lstlisting}[style=rstyle]
apps_clean |>
  filter(!is.na(Rating)) |>
  ggplot(aes(x = Type,
             y = Rating,
             fill = Type)) +
  geom_boxplot(
    width = 0.5,
    notch = TRUE,
    outlier.alpha = 0.2) +
  stat_summary(
    fun = mean,
    geom = "point",
    shape = 18, size = 3,
    color = "#E53935") +
  scale_fill_manual(
    values = c(
      Free = "#BBDEFB",
      Paid = "#FFCDD2")) +
  theme_minimal() +
  theme(legend.position =
    "none") +
  labs(title = "Rating by Type",
       y = "Rating", x = NULL)
\end{lstlisting}
    \end{column}
    \begin{column}{0.45\textwidth}
      \includegraphics[width=\textwidth]{figures_m10/box_r}
    \end{column}
  \end{columns}
\end{frame}

% ================================================================
\begin{frame}[fragile]{Boxplot --- Python (matplotlib)}
  \begin{columns}[T]
    \begin{column}{0.52\textwidth}
\begin{lstlisting}[style=pythonstyle]
free = df.loc[
    df["Type"]=="Free","Rating"]
paid = df.loc[
    df["Type"]=="Paid","Rating"]

fig, ax = plt.subplots(figsize=(5,4))
bp = ax.boxplot(
    [free, paid],
    labels=["Free", "Paid"],
    patch_artist=True,
    widths=0.45, notch=True)

bp["boxes"][0].set_facecolor(
    "#C8E6C9")
bp["boxes"][1].set_facecolor(
    "#FFCDD2")

# Overlay mean points
means = [free.mean(), paid.mean()]
ax.scatter([1, 2], means,
    c="#E53935", s=50, zorder=5,
    marker="D", edgecolors="white")

ax.set_ylabel("Rating")
ax.set_title("Rating by Type",
    fontweight="bold")
\end{lstlisting}
    \end{column}
    \begin{column}{0.45\textwidth}
      \includegraphics[width=\textwidth]{figures_m10/box_py}
    \end{column}
  \end{columns}
\end{frame}

% ================================================================
\section{Chart 5: Line Chart}
% ================================================================

\begin{frame}[fragile]{Line Chart --- R (ggplot2)}
  \begin{columns}[T]
    \begin{column}{0.52\textwidth}
\begin{lstlisting}[style=rstyle]
apps_clean |>
  mutate(year = year(
    mdy(`Last Updated`))) |>
  count(year, Type) |>
  filter(year >= 2010) |>
  ggplot(aes(x = year, y = n,
             color = Type)) +
  geom_line(linewidth = 1.2) +
  geom_point(size = 2) +
  # Direct labels at endpoint
  geom_text(
    data = . %>%
      filter(year == max(year)),
    aes(label = Type),
    nudge_x = 0.3, size = 3,
    fontface = "bold") +
  scale_color_manual(
    values = c(Free = "#1565C0",
               Paid = "#E53935")) +
  theme_minimal() +
  theme(legend.position = "none") +
  labs(title = "Apps Added/Year",
       x = "Year", y = "Count")
\end{lstlisting}
    \end{column}
    \begin{column}{0.45\textwidth}
      \includegraphics[width=\textwidth]{figures_m10/line_r}
    \end{column}
  \end{columns}
\end{frame}

% ================================================================
\begin{frame}[fragile]{Line Chart --- Python (matplotlib)}
  \begin{columns}[T]
    \begin{column}{0.52\textwidth}
\begin{lstlisting}[style=pythonstyle]
apps_clean["year"] = pd.to_datetime(
    apps_clean["Last Updated"],
    format="%B %d, %Y",
    errors="coerce").dt.year

yearly = (apps_clean
    .groupby(["year","Type"])
    .size()
    .reset_index(name="n")
    .query("year >= 2010"))

fig, ax = plt.subplots(figsize=(6,4))
for t, c in [("Free","#2E7D32"),
             ("Paid","#E53935")]:
    sub = yearly.query("Type==@t")
    ax.plot(sub["year"], sub["n"],
        "o-", color=c, lw=2, ms=5)
    ax.text(sub["year"].iloc[-1]+0.2,
        sub["n"].iloc[-1], t,
        fontsize=8, fontweight="bold",
        color=c, va="center")

ax.set_title("Apps Added per Year",
    fontweight="bold")
ax.spines["top"].set_visible(False)
ax.spines["right"].set_visible(False)
\end{lstlisting}
    \end{column}
    \begin{column}{0.45\textwidth}
      \includegraphics[width=\textwidth]{figures_m10/line_py}
    \end{column}
  \end{columns}
\end{frame}

% ================================================================
\section{Chart 6: Stacked Bar}
% ================================================================

\begin{frame}[fragile]{Stacked Bar --- R vs Python}
  \begin{columns}[T]
    \begin{column}{0.48\textwidth}
      \textbf{R (ggplot2)}
      \includegraphics[width=\textwidth]{figures_m10/stacked_r}
\begin{lstlisting}[style=rstyle]
ggplot(apps_clean,
  aes(x = `Content Rating`,
      fill = Type)) +
  geom_bar(position = "stack") +
  scale_fill_manual(values =
    c(Free="#1565C0",
      Paid="#E53935")) +
  theme_minimal()
\end{lstlisting}
    \end{column}
    \begin{column}{0.48\textwidth}
      \textbf{Python (matplotlib)}
      \includegraphics[width=\textwidth]{figures_m10/stacked_py}
\begin{lstlisting}[style=pythonstyle]
ct = pd.crosstab(
  apps_clean["Content Rating"],
  apps_clean["Type"])
ct.plot.bar(stacked=True,
  color=["#2E7D32","#E53935"],
  edgecolor="white", ax=ax)
\end{lstlisting}
    \end{column}
  \end{columns}
\end{frame}

% ================================================================
\section{Syntax Cheat Sheet}
% ================================================================

\begin{frame}{R $\leftrightarrow$ Python Complete Comparison}
  \begin{center}
    \includegraphics[width=0.92\textwidth]{figures_m10/syntax_card}
  \end{center}
\end{frame}

% ================================================================
\section{Lab Dashboard}
% ================================================================

\begin{frame}{Final Output: Six-Panel EDA Dashboard}
  \begin{center}
    \includegraphics[width=0.92\textwidth]{figures_m10/dashboard}
  \end{center}
  \begin{exampleblock}{Data Insight}
    This dashboard combines all six chart types from the lab. In practice
    you would produce this in R (\texttt{patchwork}) or Python
    (\texttt{GridSpec}) --- both are equally capable.
  \end{exampleblock}
\end{frame}

% ================================================================
\section{Synthesis}
% ================================================================

\begin{frame}{Workshop 1: What You Now Know}
  \centering
  \begin{tabular}{rlp{5cm}}
    \toprule
    \textbf{Module} & \textbf{Topic} & \textbf{Core Skill} \\
    \midrule
    M01 & Why Visualize           & Anscombe, Datasaurus \\
    M02 & Tufte's Principles      & Data-ink, lie factor \\
    M03 & Visual Perception       & Pre-attentive, Gestalt \\
    M04 & Color \& Accessibility  & Viridis, CVD, grey-accent \\
    M05 & Typography \& Layout    & Hierarchy, direct labels \\
    M06 & Design Process          & Cairo, Munzner, pipeline \\
    M07 & Critique \& Redesign    & 4-quadrant, SBI feedback \\
    M08 & R Environment           & Base R graphics \\
    M09 & Python Environment      & matplotlib OO API \\
    M10 & Comparative Lab         & R/Python parity \\
    \bottomrule
  \end{tabular}
\end{frame}

% ================================================================
\begin{frame}{Key Takeaways}
  \begin{enumerate}
    \item \textbf{R and Python produce identical results} --- the choice
          is context-dependent, not quality-dependent.
    \item The \textbf{six chart types} (bar, histogram, scatter, boxplot,
          line, stacked bar) cover the vast majority of EDA needs.
    \item Every chart applied principles from M01--M07: sorted bars (M02),
          log scales (M03), colour by type (M04), direct labels (M05),
          Tufte cleanup (M02), audience-appropriate titles (M06).
    \item \textbf{Workshop 2} builds on this foundation with the Grammar
          of Graphics: ggplot2 layers (R) and plotnine/seaborn (Python).
  \end{enumerate}
\end{frame}

% ================================================================
\begin{frame}{Workshop 1 Homework Summary}
  \begin{alertblock}{Due: Before Workshop 2}
    Combine all module homeworks into a single PDF portfolio:
    \begin{itemize}
      \item 10 module assignments (Parts A--C each)
      \item All R scripts (\texttt{.R}) and Python scripts (\texttt{.py})
      \item All exported figures (PDF vector + PNG 300 dpi)
    \end{itemize}
    Total weight: 5\% of course grade.
  \end{alertblock}

  \vspace{6pt}
  \textbf{Next}: Workshop 2 --- The Grammar of Graphics (M01: Wilkinson's Theory)
\end{frame}

\end{document}
